\section{Additive coding of the quartic source}
\label{b4:section}

The main new construction replaces the positional exponents in the
ten-integer equation of package~31 by an additive code.  The source
quartic, its complete coefficient tuple, and the outer Lucas and
relation-combining formulas are retained.  In particular, the
ten-integer theorem remains Sun's theorem \cite{Sun2026}; the result
below is an explicit improvement of its source specialization
\cite{Package31}.  Bose--Chowla codes themselves are classical
\cite{BoseChowla}.

\begin{theorem}[The explicit ten-integer instance]
\label{b4:main}
The circuit \code{ten_integer/generated/b4_signed10.json} defines a
polynomial
\[
 F_{B_4}\in\mathbb Z[f,g,h,k,l,w,x,y,z,m]
\]
of exact total degree
\[
 \boxed{\deg F_{B_4}=107\,286\,763\,576.}
\]
Its supplied division-free circuit has $2\,204$ binary additions,
subtractions, and multiplications.  With the fixed parent quartic and
compiler tuple specified in this report, ordinary mathematics
formalizable in ZFC proves
\[
 \exists(f,g,h,k,l,w,x,y,z,m)\in\mathbb Z^{10}\;(F_{B_4}=0)
 \quad\Longleftrightarrow\quad
 \neg\operatorname{Con}(\mathrm{ZFC}).
\]
Consequently absence of integer zeros is equivalent, in ZFC, to
$\operatorname{Con}(\mathrm{ZFC})$.
\end{theorem}

The proof imports the same K215 calibration, native compiler, positive
quartic reduction, and Sun bridge as the pinned construction.  It gives
a new proof of the entire source-coding interface.  Exact finite checks
supplement that proof; they are not an end-to-end formal verification of
the imported results.  Sun's ten-integer bridge is cited in its 2026
preprint version \cite{Sun2026}.  No consistency proof, globally minimal
equation, or new general ten-variable theorem is claimed.

\subsection{The fixed source and its coefficient convention}
\label{b4:source}

Let $r_1,\ldots,r_{22}$ be the fixed positive-parent residuals specified
in this report, in their literal order, and put
\[
 P_0(n_1,\ldots,n_{31})
 =\sum_{j=1}^{22}r_j(n_1+1,\ldots,n_{31}+1)^2.
\]
The inherited endpoint states
\[
 \exists\mathbf n\in\mathbb N^{31}\;P_0(\mathbf n)=0
 \quad\Longleftrightarrow\quad
 \text{ZFC has a finite syntactic refutation}.
 \tag{B4.1}\label{b4:endpoint}
\]
Here $\mathbb N$ includes zero.  The positive-to-natural shift is part
of the definition and cannot be omitted.

The four fixed coefficient ports are the complete Wang943 tuple of
package~31, with port order $(\beta,K_0,M_C,M_F)$ and bit lengths
\[
 (1064921,\ 1070004,\ 1064921,\ 1064922).
\]
The high mask $M_F$ is already shifted.  None of these ports is
existentially quantified, and no port is substituted from another
tuple.  The exact parent and compressed integer payloads are included
in \code{ten_integer/data/}.

The source has degree four and $P_0(\mathbf0)>0$.  For the latter fact,
the fifth residual $E-X(a-E)$ equals $1$ at the all-one parent
assignment.  Write
\begin{align}
 P_0(\mathbf n)&=\sum_{|\alpha|\le4}p_\alpha\mathbf n^\alpha,
 &
 c_\alpha&=(4-|\alpha|)!\prod_{i=1}^{31}\alpha_i!\,p_\alpha,
 \label{b4:coeff}\\
 H&=\max_\alpha|p_\alpha|,
 &
 \beta_S&=2^{4269882}.
 \label{b4:beta}
\end{align}
The included coefficient certificate proves
\[
 \beta_S>2\cdot24\cdot33^4H.
 \tag{B4.2}\label{b4:bound}
\]
Indeed it bounds every weighted coefficient by
$|c_\alpha|+1<2^{4269855}$, and
$2\cdot24\cdot33^4=56924208<2^{27}$.  Since the integer factorial
multiplier in \eqref{b4:coeff} is positive,
$|p_\alpha|\le|c_\alpha|$.  Thus \eqref{b4:bound} follows without
requiring decimal expansion of the enormous coefficients.

The source has $402$ nonzero weighted coefficient records.  Each is
an explicit small polynomial in the four fixed ports.  The verifier
reconstructs every record from the $22$ shifted residual squares.
The leading zero retained in each serialized exponent vector is a
historical coefficient-interface convention: it is not an input digit,
an unknown, or an extra source witness.

\subsection{What the exponent code must distinguish}
\label{b4:injectivity}

\begin{definition}
A list $0=w_0<w_1<\cdots<w_{31}$ is a quartic additive code if all
sums of four list elements, allowing repetition and disregarding
order, are distinct.
\end{definition}

Padding a multi-index with $4-|\alpha|$ copies of $w_0$ shows that this
condition is equivalent to injectivity of
\[
 \alpha\longmapsto\alpha\cdot\mathbf w
 :=\sum_{i=1}^{31}\alpha_iw_i
 \qquad (|\alpha|\le4).
 \tag{B4.3}\label{b4:injective}
\]
The positional weights $5^i$ satisfy this property.  They are not
needed for either the witness-digit mask or coefficient extraction.
An additive code gives much smaller exponents.

The code used here is
\[
\begin{gathered}
0,\ 42609,\ 49910,\ 91724,\ 155840,\ 213052,\ 290241,\ 337708,\\
340502,\ 351998,\ 388348,\ 417182,\ 480443,\ 513128,\ 534148,\\
563475,\ 579595,\ 636647,\ 639228,\ 692126,\ 725139,\ 732225,\\
748536,\ 754497,\ 817940,\ 820267,\ 829985,\ 858576,\ 870324,\\
879730,\ 907673,\ 942302.
\end{gathered}
\]
The $i$th positive entry is assigned to the $i$th coordinate of the
literal parent ordering.  Throughout this section set
\[
 T=942302,\qquad L_c=4T=3769208,\qquad M=2^{20}-1=1048575.
 \tag{B4.4}\label{b4:span}
\]

\begin{lemma}[Finite code certificate]
\label{b4:code}
The displayed list is a quartic additive code.  In fact, its four-term
multiset sums are distinct modulo $M$.
\end{lemma}

\begin{proof}
There are $\binom{35}{4}=52360$ relevant multisets.  The included
short checker enumerates all of them and verifies distinct residues
modulo $M$.  This provides a direct finite certificate independent of
the construction history.

For a mathematical construction, form the binary quotient by
$X^{20}+X^3+1$ and let $\theta$ be the residue of $X$.  Exact modular
computations give $\theta^M=1$ and
\[
 \theta^{M/p}\ne1\qquad(p=3,5,11,31,41),
\]
the complete set of prime divisors of $M$.  Hence $\theta$ has order
$M$.  Its $M$ distinct powers and zero exhaust the $2^{20}$ residues.
Every nonzero residue is therefore a unit, proving that the quotient
is a field, rather than assuming irreducibility in this step.

Let $\mathbb F_{32}$ be its $32$-element subfield and choose the
$32$ logarithms $a_c$ satisfying
\[
 \theta^{a_c}=\theta+c\qquad(c\in\mathbb F_{32}).
\]
The element $\theta$ has degree four over $\mathbb F_{32}$.  If two
sums of four logarithms agree modulo $M$, the corresponding products
of four factors $\theta+c$ agree.  Subtracting their monic
degree-four polynomials leaves a polynomial of degree at most three
over $\mathbb F_{32}$.  It must vanish identically, so unique
factorization identifies the two multisets of factors.  This is the
Bose--Chowla argument \cite{BoseChowla}.

Translate the logarithms by $-535448$ modulo $M$ and take their
representatives in $[0,M)$.  The sorted list is the displayed code.
Translation preserves equality or inequality of four-term sums
modulo $M$.  The chosen origin follows the largest cyclic gap of
the original logarithms and reduces the maximum representative to $T$.
The generator reproduces each finite-field calculation, all logarithms,
the translation, and the exhaustive multiset certificate.
\end{proof}

\subsection{The coefficient identity}
\label{b4:polarization}

Define formal polynomials
\begin{align}
 C_{\mathbf n}(t)&=\left(1+\sum_{i=1}^{31}n_it^{w_i}\right)^4,
 &
 D_c(t)&=\sum_{|\alpha|\le4}c_\alpha t^{L_c-\alpha\cdot\mathbf w}.
 \label{b4:CD}
\end{align}
Every exponent of $D_c$ is nonnegative, since
$\alpha\cdot\mathbf w\le4T=L_c$.

\begin{lemma}[Quartic coefficient recovery]
\label{b4:extract}
For every $\mathbf n$,
\[
 [t^{L_c}]C_{\mathbf n}(t)D_c(t)=24P_0(\mathbf n).
 \tag{B4.5}\label{b4:recovery}
\]
If $0\le n_i<b$ and $q_j$ denotes any coefficient of
$C_{\mathbf n}(t)D_c(t)$, then
\[
 |q_j|\le24H\left(1+\sum_i n_i\right)^4
       \le24H\bigl(1+31(b-1)\bigr)^4.
 \tag{B4.6}\label{b4:qbound}
\]
The product has degree at most $8T$.
\end{lemma}

\begin{proof}
The multinomial term indexed by $\gamma$ in $C_{\mathbf n}$ has
coefficient
\[
 \frac{4!}{(4-|\gamma|)!\prod_i\gamma_i!}\,\mathbf n^\gamma
\]
and exponent $\gamma\cdot\mathbf w$.  Its product with the
$\alpha$-term of $D_c$ reaches exponent $L_c$ exactly when
$\gamma\cdot\mathbf w=\alpha\cdot\mathbf w$.  Lemma~\ref{b4:code}
forces $\gamma=\alpha$, and the factorials cancel.  Summing proves
\eqref{b4:recovery}.

Distinct indices in $D_c$ have distinct exponents and coefficients
of magnitude at most $24H$.  The coefficients of $C_{\mathbf n}$
are nonnegative and their sum is $C_{\mathbf n}(1)$.  Bounding each
convolution coefficient by $24H C_{\mathbf n}(1)$ proves
\eqref{b4:qbound}.  Finally,
$\deg C_{\mathbf n}\le4T$ and $\deg D_c\le L_c=4T$.
\end{proof}

\subsection{The general mechanism behind the improvement}
\label{b4:general}

The coefficient step is not special to the particular quartic support.
It explains why the change can be much larger than a local circuit
simplification.

\begin{proposition}[Polynomial span at fixed degree]
\label{b4:generalspan}
Let $d\ge2$ and let $P\in\Z[x_1,\ldots,x_n]$ have total degree at
most $d$.  For every prime power $q\ge n+1$ there are fixed
integer weights
\[
 0=w_0<w_1<\cdots<w_n\le T\le q^d-2
\]
such that, on putting
\[
 c_\alpha=(d-|\alpha|)!\prod_i\alpha_i!\,p_\alpha,
 \qquad
 D_P(t)=\sum_{|\alpha|\le d}c_\alpha t^{dT-\alpha\cdot\mathbf w},
\]
one has the exact identity
\[
 [t^{dT}]
 \left(1+\sum_{i=1}^n x_it^{w_i}\right)^dD_P(t)=d!\,P(\mathbf x).
\]
In particular, choosing the least power of two $q\ge n+1$ gives
$T<[2(n+1)]^d$.
\end{proposition}

\begin{proof}
The Bose--Chowla construction in $\F_{q^d}$ gives a set of $q$
residue classes modulo $q^d-1$ with distinct unordered $d$-sums.
Indeed a primitive element has degree $d$ over $\F_q$;
equality of products of $d$ factors $\theta+c$ yields an
identically zero polynomial of degree at most $d-1$, exactly
as in Lemma~\ref{b4:code}.  Translate one class to zero, retain
it and any $n$ other classes, and take their representatives in
$[0,q^d-2]$.  Padding with zero gives the needed injectivity
for all multi-indices of total degree at most $d$.
The multinomial proof of Lemma~\ref{b4:extract}, with $d!$
in place of $24$, proves the displayed identity.  Finally
$q<2(n+1)$ for the stated choice of $q$.
\end{proof}

Thus at fixed $d$ the required exponent span is polynomial in
the number of source variables, whereas the standard positional
choice has span $(d+1)^n$.  This statement concerns the
coefficient-coding stage.  A full Diophantine reduction still
requires its bounds, masks, and both implications to be checked.
The following proof performs that work for the delivered quartic
instance.  The finite-field theorem is classical; no claim of
literature priority for this general use of additive codes is made.

\subsection{The new two-witness source interface}
\label{b4:interface}

For integer variables $f,g$, define the following expressions in the
displayed order:
\begin{align}
 b&=1+3f,&
 \mathcal B&=\beta_Sb^4,&
 \ell&=\sum_{i=1}^{31}\mathcal B^{w_i},
 \label{b4:source1}\\
 D_c&=D_c(\mathcal B),&
 N_0&=\mathcal B^{T+1},&
 N_1&=4\mathcal B^{8T+1},
 \label{b4:source2}\\
 N&=N_0N_1,&
 \mathscr C&=b\mathcal B^T,&
 M_c&=N_0-1-(b-1)\ell,
 \label{b4:source3}\\
 T_c&=M_c+(\mathcal B-2)\mathcal B^{L_c}N_0,
 \label{b4:source4}\\
 J_c&=2(1+g)^4D_c+\sum_{j=0}^{8T}\mathcal B^{j+1},
 \label{b4:source5}\\
 S_c&=g+J_cN_0,&
 R_c&=(S_c+T_c+1)N+T_c+1,
 \label{b4:source6}\\
 X&=(N-1)R_c,&
 Y&=N^2.
 \label{b4:source7}
\end{align}
All exponents and sum limits are fixed integers.  These expressions
are ordinary polynomials in $f,g$ with fixed integer coefficients.

\begin{theorem}[Complete source interface]
\label{b4:sourceinterface}
The following implications hold.
\begin{enumerate}
\item If $P_0$ has a natural zero, then for every positive bound $Z$
there exist $f\ge Z$ and $b\le g<\mathscr C$ such that $b$ is a
square power of two and
\[
 Y\mid\binom{2X}{X}.
 \tag{B4.7}\label{b4:binomial}
\]
\item If $f\ne0$, $0\le g<2\mathscr C$, $b$ is a square power of two,
and \eqref{b4:binomial} holds, then $P_0$ has a natural zero.
\item Already under $f\ne0$, $b$ square, and
$0\le g<2\mathscr C$, without a power-of-two assumption on $b$,
\[
 3\mid X,\qquad X>3b,\qquad Y>\max\{b,256\}.
 \tag{B4.8}\label{b4:bridgebounds}
\]
\end{enumerate}
Moreover $\mathscr C=2^{4269882T}b^{4T+1}$, with an even positive
exponent of two, so the bound has exactly the form required by
Sun's source interface.
\end{theorem}

\begin{proof}
We first establish signs and bounds before assuming a source zero or
binomial divisibility.  Since $b=1+3f$ is a square and $f\ne0$,
one has $f>0$ and $b\ge4$.  The coefficient bound
\eqref{b4:bound} implies
\begin{align}
 \frac{\mathcal B}{2}
 &>24H\bigl(1+31(b-1)\bigr)^4,
 \label{b4:firstbound}\\
 \mathcal B
 &>2+2\cdot24H(32b)^4.
 \label{b4:secondbound}
\end{align}
For the second inequality, replacing $33$ by $32$ in the bound for
$\beta_S$ leaves more than two units of margin after multiplication
by $b^4$, because $H\ge1$ and $b\ge4$.

The constant coefficient $P_0(\mathbf0)$ is positive.  Thus $D_c>0$:
its leading term is $24P_0(\mathbf0)\mathcal B^{L_c}$ and the
absolute sum of all lower terms is at most
\[
 24H\sum_{j=0}^{L_c-1}\mathcal B^j<\mathcal B^{L_c}.
\]
We also have
$D_c\le24H\sum_{j=0}^{L_c}\mathcal B^j$.
The base-$\mathcal B$ digits of $M_c$ are $\mathcal B-b$ at
positions $w_i$ and $\mathcal B-1$ at the other positions from
$0$ through $T$.  Therefore
\[
 0\le M_c<N_0,\qquad
 0<(\mathcal B-2)\mathcal B^{L_c}<N_1,\qquad
 0\le T_c<N.
 \tag{B4.9}\label{b4:blocksT}
\]

For $0\le g<2\mathscr C$, the inequalities
$2b\le\mathcal B$ and $g\in\mathbb Z$ give
$g<N_0$ and $1+g\le2b\mathcal B^T\le32b\mathcal B^T$.
Using \eqref{b4:secondbound} and the bound for $D_c$,
\begin{align*}
 0<J_c
 &<
 \frac{(\mathcal B-2)+\mathcal B}{\mathcal B-1}
       \mathcal B^{8T+1}\\
 &=2\mathcal B^{8T+1}<N_1.
\end{align*}
Consequently $0\le S_c<N$.  These are unconditional bounds on
the stated source range.

Because $\beta_S$ is an even power of two and $b\equiv1\pmod3$,
we have $\mathcal B\equiv1\pmod3$ and
$N=4\mathcal B^{9T+2}\equiv1\pmod3$.  Hence $3\mid X$.
Also $R_c\ge N+1$, $N>\mathcal B\ge b$, and $N\ge4$, so
$X=(N-1)R_c>3b$ and $Y=N^2>\max\{b,256\}$.
This proves \eqref{b4:bridgebounds}.

Now assume additionally that $b$ is a power of two.  Then
$\mathcal B,N_0,N_1,N$ are powers of two.  For nonnegative integers
$a,d$, write $\tau_2(a,d)=0$ when their binary addition has no
carries, equivalently when their bitwise AND is zero.  The mask
description gives
\begin{multline}
 0\le g<2\mathscr C,\quad \tau_2(g,M_c)=0\\
 \Longleftrightarrow\quad
 g=\sum_{i=1}^{31}n_i\mathcal B^{w_i}
 \quad\hbox{for some }0\le n_i<b.
 \tag{B4.10}\label{b4:mask}
\end{multline}
Indeed the digit $\mathcal B-b$ has ones exactly in the high binary
positions of one $\mathcal B$-digit, whereas $\mathcal B-1$ has
ones in every position.  The only free bits of $g$ are therefore
the low $\log_2b$ bits at the distinct slots $w_i$.  The bound
$g<N_0$ excludes all higher slots.  Conversely every displayed
code has the required bit pattern and is less than $\mathscr C$.
No superincreasing property of the weights is used.

For decoded $g$, write
\[
 (1+g)^4D_c
 =\sum_{j=0}^{8T}q_j\mathcal B^j
\]
using the formal coefficients of Lemma~\ref{b4:extract}, and put
\[
 J'=\sum_{j=0}^{8T}
          \left(q_j+\frac{\mathcal B}{2}\right)\mathcal B^j.
\]
Equations \eqref{b4:qbound} and \eqref{b4:firstbound} show that
every displayed digit of $J'$ is an integer strictly between
$0$ and $\mathcal B$.  In particular this is its actual
base-$\mathcal B$ expansion, even when some $q_j$ are negative.
Since $J_c=2J'$,
\[
 \tau_2(J_c,(\mathcal B-2)\mathcal B^{L_c})=0
 \ \Longleftrightarrow\
 \tau_2\!\left(J',
       \left(\frac{\mathcal B}{2}-1\right)\mathcal B^{L_c}\right)=0.
\]
The mask on the right occupies precisely the low bits of the
digit at $L_c$.  A digit in $[1,\mathcal B)$ avoids it exactly
when it equals $\mathcal B/2$.  By \eqref{b4:recovery}, this is
equivalent to $q_{L_c}=24P_0(\mathbf n)=0$.

The two block decompositions, together with $S_c,T_c<N$, give
\[
 \tau_2(S_c,T_c)=0
 \ \Longleftrightarrow\
 \tau_2(g,M_c)=0
 \ \land\
 \tau_2(J_c,(\mathcal B-2)\mathcal B^{L_c})=0.
 \tag{B4.11}\label{b4:blocksplit}
\]
Sun's binary two-block lemma \cite[Lemma~2.2]{Sun2021} states,
for $0\le S_c,T_c<N$ with $N$ a power of two and the exact
$R_c,X$ above, that
\[
 \tau_2(S_c,T_c)=0
 \quad\Longleftrightarrow\quad
 N^2\mid\binom{2X}{X}.
 \tag{B4.12}\label{b4:Kummer}
\]
Its hypotheses have all been established before invoking it.
Combining \eqref{b4:mask}--\eqref{b4:Kummer} proves soundness.

For completeness, suppose $P_0(\mathbf n)=0$.  Choose $s$ so large
that $b=2^{2s}>\max_i n_i$ and $(b-1)/3\ge Z$.  Since
$4^s\equiv1\pmod3$, $f=(b-1)/3$ is an integer.  The source
zero is not the all-zero vector, so
$g=\sum_i n_i\mathcal B^{w_i}\ge\mathcal B\ge b$.
The usual geometric-sum estimate gives $g<b\mathcal B^T=\mathscr C$.
The mask and coefficient identities yield
\eqref{b4:binomial}.  Finally the asserted monomial formula for
$\mathscr C$ follows directly from $\mathcal B=\beta_Sb^4$.
\end{proof}

\subsection{The complete outer polynomial}
\label{b4:outer}

We now give the full outer formula, so that the artifact is specified
by fixed constants, finite sums, and ordinary polynomial operations.
Only the ten lowercase coordinates listed in
Theorem~\ref{b4:main} are unknowns.  Every other symbol is a circuit
register.

Starting with the new $b,X,Y,\mathscr C$ above, set
\begin{align}
 L_s&=lY,& U&=2L_sX,\notag\\
 V&=4gwY,& W_s&=bw,
 \label{b4:bridge1}\\
 A&=U(V+1),& B_s&=2X+1,\notag\\
 C&=B_s+(A-2)h,& P_n&=U^2V,
 \label{b4:bridge2}\\
 K&=X+1+k(P_n-2),&
 S_s&=2A-5,
 \label{b4:bridge3}\\
 T_s&=3W_sC-2(W_s^2-1),&
 R_2&=(P_n^2-4)K^2+4,
 \label{b4:bridge4}
\end{align}
\begin{equation}
 \Delta=8\mathscr C^3gK^2
       -g^2\{32(C-KL_s)^2\mathscr C^3+g^2K^2\},
 \label{b4:bridge5}
\end{equation}
\begin{align}
 O&=f^2l^2\Delta,&
 \xi&=zx.
 \label{b4:bridge6}
\end{align}
The modified Lucas certificate is
\begin{align}
 D_s&=(A^2-4)C^2+4,&
 E&=C^2D_s\xi,
 \label{b4:lucas1}\\
 F_s&=4(A^2-4)E^2+1,&
 G_*&=2(A+2)(A-2)^2E^2-1,
 \label{b4:lucas2}\\
 \Omega&=16E,&
 H_*&=C+(B_s-C)F_s+\Omega F_sy,
 \label{b4:lucas3}\\
 I_*&=(G_*^2-1)H_*^2+1.
 \label{b4:lucas4}
\end{align}
These are the same bridge and Lucas formulas as in package~31
\cite{Package31}, using Sun's modified certificate and interpolation
theorem \cite{Sun2026}.

For the divisibility and nonvanishing merger, define
\begin{align}
 r&=T_s+S_sy,&
 \Sigma&=\xi S_s^5,
 \label{b4:merge1}\\
 \Theta&=(r^3-rS_s^2-S_s^3)(r^2+23S_s^2).
 \label{b4:merge2}
\end{align}
The three radicands and their separating weights are
\begin{align}
 r_0&=b\{1+(3O-4)b^2z^2\},&
 r_1&=R_2,&
 r_2&=D_sF_sI_*,
 \label{b4:radicands}\\
 \omega_0&=1,&
 \omega_1&=r_0,&
 \omega_2&=r_0(P_n^2+K^2+3).
 \label{b4:weights}
\end{align}
Finally put
\begin{align}
 Z_c&=\Sigma m-\Theta,&
 A_i&=\Sigma^2\omega_i^2r_i\quad(0\le i\le2),
 \label{b4:norm1}\\
 V_n&=Z_c^2+A_2-A_0-A_1,
 \label{b4:norm2}\\
 W_n&=V_n^2+4Z_c^2A_2-4A_0A_1,
 \label{b4:norm3}\\
 \boxed{F_{B_4}}&\boxed{{}=W_n^2-16V_n^2Z_c^2A_2.}
 \label{b4:norm4}
\end{align}
The requested equation is $F_{B_4}=0$.  No radical in the proof below
is part of its definition.  There are no divisions or variable
exponents in \eqref{b4:source1}--\eqref{b4:norm4}.

\subsection{Why the finalizer retains both implications}
\label{b4:finalizer}

For clarity we record the inherited finalizer arguments, with their
attribution to package~31 \cite{Package31}.  They ensure that the new
source interface is being composed with the precise outer hypotheses.

\begin{lemma}[The quintic merger]
\label{b4:quintic}
The polynomial
$\pi(v)=(v^3-v-1)(v^2+23)$ has a root modulo every positive integer
and no rational root.  If $S_s\ne0$ and $\xi\ne0$, then
$\Sigma\mid\Theta$ implies $S_s\mid T_s$.
Conversely, when $S_s\mid T_s$, a residue class of $y$ modulo
$|\xi|$ makes $\Sigma\mid\Theta$.  For every integer assignment
with $S_s\ne0$, $\Theta\ne0$.
\end{lemma}

\begin{proof}
The cubic has discriminant $-23$.  At $p=23$ it has the simple
root $3$.  At an odd prime $p\ne23$, either it has a simple root
or it is irreducible.  In the latter case Frobenius is an even
three-cycle, so the discriminant is a square modulo $p$ and
$v^2+23$ has a simple root.  Hensel lifting handles all odd prime
powers.  Since $-23\equiv1\pmod8$, the quadratic has a root in
the $2$-adic integers as well.  The Chinese remainder theorem
gives a root modulo every positive integer.

The rational-root theorem excludes rational roots of the cubic,
whose only possible integer roots are $\pm1$; the quadratic is
positive on $\mathbb Q$.  Consequently
$\Theta=S_s^5\pi((T_s+S_sy)/S_s)$ cannot vanish.

If $S_s\mid T_s$, write $T_s=S_sq$ and choose
$y\equiv c-q\pmod{|\xi|}$ for a root $c$ of $\pi$ modulo $|\xi|$.
Then $\Theta=S_s^5\pi(q+y)$ is divisible by $\Sigma$.
Conversely suppose that a prime satisfies
$a=v_p(S_s)>d=v_p(T_s)$.  Then $v_p(r)=d$; the leading term of
the homogeneous cubic factor has uniquely smallest valuation
$3d$, and that of the quadratic factor has uniquely smallest
valuation $2d$.  Thus $v_p(\Theta)=5d<5a$, contradicting
$S_s^5\mid\Theta$.  Therefore $S_s\mid T_s$.
\end{proof}

\begin{lemma}[Positivity and separating weights]
\label{b4:separation}
For the displayed polynomial expressions, if $z\ne0$ then
$|r_0|>1$ and
\[
 \sqrt{|r_1|}\le P_n^2+K^2+2.
 \tag{B4.13}\label{b4:r1bound}
\]
Any rational signed sum
$\pm\sqrt{r_0}\pm\omega_1\sqrt{r_1}\pm\omega_2\sqrt{r_2}$
forces every $r_i$ to be an integer square.
Moreover, $r_0$ square and $z\ne0$ imply that $b$ is a positive
square and $O>0$.
\end{lemma}

\begin{proof}
The congruence $b\equiv1\pmod3$ excludes $b=0,-1$.  If
$r_0=\pm1$, its factor $b$ must be $1$, hence $f=O=0$ and
$r_0=1-4z^2$, a contradiction.  If $r_0=0$, its bracket vanishes,
forcing $|b|=|z|=1$ and $3O-4=-1$; again $b=1$ gives $O=0$,
a contradiction.  This proves $|r_0|>1$.

For \eqref{b4:r1bound}, use
$|r_1|\le(P_n^2+4)K^2+4$.  Subtracting this bound from
$(P_n^2+K^2+2)^2$ leaves
$P_n^4+K^4+P_n^2K^2+4P_n^2\ge0$.
It follows that
\[
 |\omega_1|>\sqrt{|r_0|},\qquad
 |\omega_2|>\sqrt{|r_0|}
                    +|\omega_1|\sqrt{|r_1|}.
 \tag{B4.14}\label{b4:strictweights}
\]
In the multiquadratic splitting field, any automorphism changing a
nonzero term of a rational signed sum has a largest changed index.
The magnitude of that term is at least the magnitude of its weight,
since its nonzero radicand is an integer.  By
\eqref{b4:strictweights}, all earlier terms together are too small
to cancel its change.  This contradiction shows that every root
term is rational.  A rational square root of an integer is integral.
The argument also covers negative candidate radicands by complex
absolute values, and zero radicands require no recovery.

Finally the two factors of $r_0$ are coprime, since the bracket
is $1$ modulo $b$.  If their product is a nonzero square, each
factor is a signed square with the same sign.  A negative square
cannot be $1$ modulo $3$, so $b$ is a positive square and its
bracket is a square.  If $O\le0$ and $bz\ne0$, that bracket
would be at most $1-4=-3$.  Hence $O>0$.
\end{proof}

\begin{proof}[Proof of Theorem~\ref{b4:main}: equivalence]
The polynomial identity underlying the finalizer is
\[
 F_{B_4}
 =\prod_{\varepsilon\in\{-1,1\}^3}
   \left(Z_c+\sum_{i=0}^2\varepsilon_i\sqrt{A_i}\right).
 \tag{B4.15}\label{b4:normidentity}
\]
It follows by successively multiplying conjugate factors, or by
expanding the $35$ terms of the three-radicand norm.  The delivered
checker proves this exact formal identity independently.

Suppose $F_{B_4}=0$.  The integer $S_s=2A-5$ is odd and nonzero,
so Lemma~\ref{b4:quintic} gives $\Theta\ne0$.  If $\Sigma=0$,
then every $A_i=0$ and \eqref{b4:norm4} equals $\Theta^8\ne0$.
Thus $\Sigma\ne0$, and $\xi=zx\ne0$.  Divide a vanishing factor
of \eqref{b4:normidentity} by $\Sigma$.  Lemma~\ref{b4:separation}
makes every $r_i$ a square; the resulting signed sum of roots is
an integer, so $\Sigma\mid\Theta$.
Lemma~\ref{b4:quintic} gives $S_s\mid T_s$, and
Lemma~\ref{b4:separation} gives $b$ square and $O>0$.
These, $\xi\ne0$, $R_2$ square, and $D_sF_sI_*$ square are
exactly the hypotheses of Sun's reverse bridge
\cite[Proposition~6.2]{Sun2026}.  It yields the source range,
the power-of-two condition, and binomial divisibility required
by Theorem~\ref{b4:sourceinterface}.  Hence $P_0$ has a natural
zero, and \eqref{b4:endpoint} yields a finite ZFC refutation.

Conversely suppose $P_0$ has a natural zero.
Theorem~\ref{b4:sourceinterface} supplies arbitrarily large $f$;
choose $f\ge2$.  Sun's positive bridge
\cite[Proposition~6.1]{Sun2026} supplies the other initial
witnesses with $b$ square, $R_2$ square, $S_s\mid T_s$,
$C=u_{B_s}(A)$, and $O>0$.  As
$O=f^2l^2\Delta$ is a positive integer, $O\ge4$.
The positive nonsquare integer $d=3O-4\equiv2\pmod3$
has a Pell solution with its second coordinate a positive
multiple of $b$: take a positive Pell unit to a power equal
to the order of its class modulo $b$.  Write that second
coordinate as $bz$ with $z>0$.  Then $r_0$ is a square.

Sun's full residue-class interpolation theorem
\cite[Theorem~4.1]{Sun2026} supplies a positive $\xi$ divisible
by this $z$, and allows $y$ in the residue class required by
Lemma~\ref{b4:quintic}, while making $D_sF_sI_*$ a square.
Set $x=\xi/z$.  Thus all $r_i$ are squares,
$\Sigma\ne0$, and $\Sigma\mid\Theta$.  Choose the integer $m$
so that a selected signed factor of \eqref{b4:normidentity}
vanishes.  This gives $F_{B_4}=0$.

All new steps have been proved above.  The substantial imported
machine/compiler and Lucas results are explicitly identified;
their ordinary arithmetic proofs, not the soundness of any
proof assistant, supply the asserted ZFC-provability.
\end{proof}

\subsection{Circuit identity, degree, and reproducibility}
\label{b4:verification}

The sparse coefficient evaluator uses the division-free identity
\[
 \mathcal B^{4T-\alpha\cdot\mathbf w}
 =(\mathcal B^T)^{4-|\alpha|}
       \prod_{i=1}^{31}
          (\mathcal B^{T-w_i})^{\alpha_i}.
 \tag{B4.16}\label{b4:horner}
\]
All exponents are nonnegative.  A multivariate Horner scheme shares
the powers and coefficient recipes.  Equation~\eqref{b4:horner}
is valid for arbitrary integer assignments, including
$\mathcal B=0$; no division by $\mathcal B$ is concealed.
The geometric sum in \eqref{b4:source5} uses binary identities
for fixed powers and fixed-length sums.

The main circuit contains $485$ additions, $25$ subtractions, and
$1694$ multiplications, for $2204$ operations.  There are $79$
fixed coefficient leaves after reuse and literal simplification.
The general source-degree bounds in this instance are
\[
 \deg X=108T+24,\qquad
 \deg Y=72T+16,\qquad
 \deg\mathscr C=4T+1.
 \tag{B4.17}\label{b4:sourcedegrees}
\]
Propagating degree bounds through the circuit yields
\[
 \deg F_{B_4}\le113856T+27064=107286763576.
 \tag{B4.18}\label{b4:degreeupper}
\]
To certify attainment rather than just an upper bound, substitute
one indeterminate for all ten unknowns.  With the actual fixed
coefficient payloads, the coefficient at the bound in
\eqref{b4:degreeupper} is $16$ modulo $101$.  It is also $599$
modulo $1009$ and $622356$ modulo $1000003$.  A single nonzero
such coefficient proves attainment, so the stated total degree
is exact.  The checker performs exact leading-coefficient
arithmetic modulo the primes; this is not inference from
random numerical agreement.

For comparison, package~31's positional construction has exact degree
\[
 126496\cdot5^{31}+27064
 =589042901992797851562527064
\]
and $1442$ operations.  The new degree is smaller by a factor of
approximately $5.49\times10^{15}$, at the cost of $762$ additional
arithmetic operations.  Thus the degree improvement is substantive,
while the operation-count tradeoff remains explicit.

The self-contained commands are
\begin{quote}
\code{python3 b4_code.py}\\
\code{python3 build_b4_sun10.py}\\
\code{python3 verify_b4_sun10.py}
\end{quote}
from \code{ten_integer/}.  They use only the Python standard
library.  In addition to the degree proof, they regenerate the
finite field and all $52360$ code sums, reconstruct all $402$
weighted coefficient records, check the fixed coefficient payloads
and bounds, and expand the norm identity exactly.  A separate
scalar transcription evaluates the literal $402$-term coefficient
sum and all named registers modulo three primes, independently
of the Horner circuit.  Those modular transcription checks are
regression evidence; the preceding lemmas supply the universal
equivalence.

The archive also includes four finalizer variants using more generic
weights or the literal Sun four-square merger.  They share the same
source adapter and fixed tuple.  A clean isolated rebuild reproduces
every generated JSON file byte for byte.  Source pins, coefficient
bindings, licenses, and inherited attribution notices accompany
the construction.
